\subsection{Identified failures in the full-geometry checks}
The larger geometry checks expose limitations that the preceding component
suite did not resolve. These failures concern numerical paths in specified
velocity models; they are not evidence for excluding any station or event.
All identified inputs and failed results remain in the experiment.

The completed third Hayward preflight initializes all 14,262 absolute rays,
but 139 fail the local stopping rule at 0.5 km field and path spacing.
These failures span 128 events and 80 channel epochs. Of the failed rays,
121 exhaust 1,000 updates; the remaining 18 stop earlier. All 32 rays in
the separate refinement sample and all 24 endpoint perturbations meet their
local stopping rules. Passing this sample therefore does not establish
convergence on the full observation support. The terminal driver records
a failed preflight and does not start an inversion. Its complete identified
tables, configuration, source snapshot, and failure review are preserved in
\path{data/development/hayward_first_arrival_003_review/}.

Continuous layered-ray enumeration has also completed for all 14,262
initial rays. Every reference is complete and valid, with no near ties;
9,069 first arrivals are direct and 5,193 turn below the source. Production
times exceed the reference by a median 37.803 microseconds and a maximum
2.232 ms. Only 603 rays satisfy both existing analytic-control limits.
The largest time and source-vector discrepancies belong to a locally
converged ray: event \texttt{NC/30204761} to \texttt{NC.JJR..EHZ},
absolute row 14099. Its source-vector error norm is 0.033571 s/km.
This outlier requires an identified branch/path investigation independent
of the 139 failed stopping flags. All reference candidates, original flags,
identities, and frozen inputs are in
\path{data/development/hayward_first_arrival_003_1d_support/}.

Event \path{quakeml:NC.anss.org/Event/NC/390228} and channel
\texttt{NC.CCY..EHZ} define Hayward absolute observation row 9291.
At 0.5 km field and initial path spacing,
the selected FMM-started path rejects its line search after 208 accepted
updates. Its time is 3.645674501 s and projected gradient is
$2.7990\times10^{-3}$ s/km. Increasing the iteration budget from 1,000 to
4,000 leaves this result unchanged. Observation hooks reproduce all eight
uninstrumented candidate calculations exactly, including the paths,
derivatives, iteration counts, and original stopping flags.

Two adjacent knots approach the 4 km depth interface within
$2.18\times10^{-13}$ km, leaving an approximately 1.01 km, nearly horizontal
segment. Inserting a free midpoint preserves the exact curve and action.
The new midpoint's depth derivative is $-0.001645449$ s/km; moving it
0.001 km deeper lowers time by $1.28394\times10^{-6}$ s. The velocity
increase below the interface reduces slowness, whereas a small depth tent
has zero first-order length cost in the horizontal limit. Thus the retained
coarse edge omits a feasible descent variation. The calculation uses the
saved coordinates without moving them onto the interface.

Across all 38 segment insertions and 190 tested perturbations, the exact
depth-only action changes by at most $3.04\times10^{-18}$ s on insertion,
and the production quadrature agrees with the exact action within
$1.06\times10^{-15}$ s. This excludes an action-integration discrepancy
at the tested scale. At 0.25 km field and initial path spacing, the
FMM-started path meets the unchanged combined stopping rule in 41 updates,
with time 3.645534035 s and gradient $5.46\times10^{-6}$ s/km. This resolves
the local stopping failure for this refined calculation. Independent
continuous layered-ray enumeration finds a unique first arrival at
3.645479994 s, turning at depth 4.000928869 km. The coarse selected path
is therefore late by 194.506 microseconds; the refined FMM path remains
late by 54.041 microseconds, with maximum source-derivative error
$1.51\times10^{-5}$ s/km. Its local flag does not establish continuous-model
accuracy. The records
are under \path{data/development/initial_ray_first_arrival_003_9291_audit001/}
and \path{initial_ray_first_arrival_003_9291_midpoint001/} in the same parent;
the independent reference is in
\path{initial_ray_first_arrival_003_9291_reference001/}.

A further test retains the 0.5 km initialization field and refines only
the path. At 0.125 km spacing, the selected FMM-started path stops after
20 updates, with time error 13.075 microseconds and source-vector error
norm $2.66\times10^{-6}$ s/km. At 0.0625 km, these errors fall to
3.267 microseconds and $7.88\times10^{-7}$ s/km. Both selected results meet
the existing analytic-control limits of 25 microseconds and
$3\times10^{-6}$ s/km. The 0.125 km straight-started candidate exceeds the
source-derivative limit despite local convergence. All four candidates are
retained in \path{initial_ray_first_arrival_003_9291_fine_path001/} under
the same development-data directory. This single-ray result motivates
further checks; it does not validate a full-cohort discretization.

Path refinement resolves both stopping and accuracy failures on the
139-ray subset (Table~\ref{tab:failed-ray-refinement}). Continuous layered-ray
enumeration finds a unique first arrival for every ray. At 0.125 km spacing,
all time errors satisfy the 25 microsecond control limit, but 16 source-vector
errors exceed $3\times10^{-6}$ s/km. Three source errors remain above that
limit at 0.0625 km. All 139 rays meet both limits at 0.03125 km, with maximum
time error 0.818 microseconds and source-vector error $1.83\times10^{-6}$
s/km. These calculations retain the frozen executable and original
1,000-update budget. They validate this subset at the initial model, not
the full observation support or a later three-dimensional model. The
four identified variants and their references are saved under
\path{data/development/hayward_initial_failure_refinement_001/} through
\path{hayward_initial_failure_refinement_004_reference/}.

\begin{table}[ht]
\centering
\small
\begin{tabular}{rrrrrr}
\toprule
Path (km) & Field (km) & Local passes & Both limits & Max time ($\mu$s) & Max source (s/km)\\
\midrule
0.25 & 0.25 & 134 & 114 & 54.041 & $9.80\times10^{-5}$\\
0.125 & 0.5 & 139 & 123 & 13.075 & $2.87\times10^{-5}$\\
0.0625 & 0.5 & 139 & 136 & 3.267 & $4.92\times10^{-6}$\\
0.03125 & 0.5 & 139 & 139 & 0.818 & $1.83\times10^{-6}$\\
\bottomrule
\end{tabular}
\caption{Refinement of all 139 rays that failed the third Hayward preflight.
Counts are out of 139. Accuracy uses absolute travel-time error and the
Euclidean norm of the three-component source-derivative error relative to
the continuous initial-model reference. The limits remain 25 microseconds
and $3\times10^{-6}$ s/km. The first row refines both discretizations;
the remaining rows retain the same 0.5 km initialization field.}
\label{tab:failed-ray-refinement}
\end{table}

The Loma Prieta check also separates local convergence from arrival-branch
selection. The terminal record retains all 104 requested samples.
For 38 rays, an independently
integrated path that grazes the maximum-speed plane $x=0$ is admissible.
Twenty-three production and refined actions exceed that candidate. Three
production paths report local convergence while remaining 25.693, 20.765,
and 37.159 ms slower:
event \texttt{NC/20091683} to \texttt{NC.HLT..EHZ} and
\texttt{NC.HQR..EHZ}, and event \texttt{NC/256583} to
\texttt{NC.JSC..EHZ}, respectively. At refined spacing their excesses fall
to 0.166, 0.187, and 0.133 ms, but none meets the local stopping rule.
The explicit candidate is an upper bound on the first
arrival, not a complete enumeration of branches. Its absence for another
ray does not justify excluding that ray. Full native identities, epochs,
selector flags, and captured rows are retained in
\path{data/development/loma_geometry_preflight/all_guides_002/}.

The completed Loma Prieta diagnostic retains all 104 requested rays.
Initializers succeed for every ray at both spacings, but 31 rays fail the
local stopping rule at each spacing. The largest time change on refinement
is 37.026 ms. Endpoint finite-difference comparisons reach approximately
$1.26\times10^{-3}$ s/km, but only 13 of 24 positive perturbations and
12 of 24 negative perturbations converge locally. These discrepancies
cannot be interpreted as validated derivative-error estimates. No
full synthetic observation set or Loma Prieta inversion follows this check.
The identified terminal tables are in \path{data/loma_lom_001/}; the review
and execution provenance are in
\path{data/development/loma_geometry_preflight/001/}.

A separate exact control identifies a defect in using an ordinary gradient
at an interface. For $v(x)=6$ km/s on $x\leq0$ and $v(x)=6-2x$ km/s on
$0\leq x\leq1$ km, a straight 2 km path on $x=0$ attains the lower bound
$T=1/3$ s. The current trace returns this action but reports nonconvergence,
using a midpoint normal gradient of $1/18$ s/km from the right-hand cell.
The true action has unequal one-sided derivatives there. The returned
source-normal partial is $1/36$ s/km, whereas the optimized continuous
arrival has zero source-normal derivative at this control. Relabeling its
stopping flag would therefore leave the sensitivity problem unresolved.
The reproducible control is in \path{data/development/interface_stationarity/}.
This component defect is distinct from the off-interface, slowly optimized
Loma paths already examined. No convergence-only correction has been
applied, and these checks have not produced a completed corrected inversion.
